% ═══ Results: Complete Circuit (Positive + Flipped + Conservation) ═══
\section{The Complete Circuit Combines Positive, Flipped, and Conservation Channels}
\label{sec:res-complete}

We constructed a two-sided classifier over a shared feature space of
physicochemical groups, separation bins, and conservation-pair types
(10 bits, 1024 cells).  Two ternary tables were labelled from identical pooled
counts: the positive table marks cells whose contact rate significantly exceeds
the base rate, while the flipped table marks cells whose rate is significantly
below it.  A pair receives a definitive CONTACT call if its cell is ON in the
positive table and a definitive FORBIDDEN call if its cell is ON in the
flipped table; all other pairs are abstained.

On held-out proteins, the positive channel achieved precision 0.092
($2.6\times$ base) at recall 0.203.  The flipped channel achieved specificity
of \textbf{98.3\%}: forbidden-called cells showed exactly the predicted
$\leq$1.75\% contact rate on unseen proteins, confirming that negative-selection
rules generalise across protein families.  Combined definitive-call MCC was
$+0.169$ at 22\% coverage.

Pairs in which both alignment columns are fully conserved (entropy $< 0.001$)
showed a contact rate of 11.4\%, representing a $3.2\times$ enrichment.
However, a finer analysis using continuous minimum-pair entropy revealed that
absolutely invariant positions are depleted of contacts ($0.69\times$), while
near-conserved positions (entropy $0.001$--$0.05$) are enriched ($1.35\times$).
This non-monotonicity indicates that structural-core membership requires
residues to be constrained but not absolutely frozen; catalytic or functional
constraints can lock a position without placing it in a contact-rich
environment.
